\chapter{Polynomial Interpolation}
\label{ch:ii20}

Polynomial interpolation reconstructs a polynomial from finite data. Lagrange and Newton forms expose uniqueness, computation, divided differences, and the remainder term.

\section*{Learning goals}
After this chapter, the reader should be able to:
\begin{itemize}
\item construct Lagrange and Newton interpolating polynomials from given data;
\item prove uniqueness of the degree-at-most-\(n\) interpolant from the root-counting argument;
\item compute divided differences and update a Newton interpolant when a new node is added;
\item derive or apply the interpolation remainder formula to bound the error;
\item compare equally spaced and Chebyshev-type node choices through the node polynomial;
\item compute or estimate the Lebesgue function $\lambda_n(x)=\sum_j|L_j(x)|$ and interpret the Lebesgue constant as the operator norm controlling interpolation sensitivity;
\item explain the Runge phenomenon as a stability failure of high-degree equally spaced interpolation rather than a contradiction of exact interpolation or polynomial density;
\item translate interpolation into a Vandermonde linear system and explain why distinct nodes are essential;
\end{itemize}
\section*{Conceptual roadmap}
\[
\boxed{\text{definition}\;\longrightarrow\;\text{estimate}\;\longrightarrow\;\text{theorem}\;\longrightarrow\;\text{application}.}
\]

\section{Interpolation problem}
Given distinct nodes and values, seek a polynomial matching all data.

\section{Lagrange basis}
Cardinal polynomials isolate one node at a time.

\section{Uniqueness}
A polynomial of degree at most $n$ matching $n+1$ distinct data points is unique.

\section{Newton form}
Divided differences produce a nested incremental representation.


% BEGIN VOL02-EXPANSION II20-example-01
\begin{example}[Newton form updates without rebuilding the interpolant]\label{ex:ii20-expand-01}
For data
\[
(0,1),\quad(1,2),\quad(2,5),
\]
the divided differences are
\[
f[0]=1,\qquad f[0,1]=1,\qquad f[1,2]=3,\qquad
f[0,1,2]=1.
\]
Hence
\[
p_2(x)=1+x+x(x-1)=1+x^2.
\]
If a new node is added, Newton form appends one new divided-difference term
rather than rebuilding the whole polynomial. Evaluating at the three old nodes
checks the interpolation conditions.
\end{example}
% END VOL02-EXPANSION II20-example-01
\section{Divided differences}
Recursive coefficients encode discrete derivatives.

\section{Error formula}
The interpolation remainder involves the next derivative and the node polynomial.


% BEGIN VOL02-EXPANSION II20-example-02
\begin{example}[Interpolation remainder becomes an error bound]\label{ex:ii20-expand-02}
Interpolate \(f(x)=e^x\) at \(0\) and \(1\) by a line \(p_1\). For
\(x\in[0,1]\),
\[
f(x)-p_1(x)=\frac{e^\xi}{2}x(x-1)
\]
for some \(\xi\in(0,1)\). Since \(e^\xi\le e\) and
\[
|x(x-1)|\le\frac14,
\]
we get
\[
|f(x)-p_1(x)|\le\frac{e}{8}.
\]
The remainder formula separates smoothness control from node geometry.
\end{example}
% END VOL02-EXPANSION II20-example-02
\section{Node choice}
% VOL02-DEPENDENCY-PREVIEW-II20-CHEBYSHEV
Equally spaced nodes can be unstable at high degree; special nodes reduce worst-case error. Chebyshev-type nodes are mentioned here as motivation and are developed systematically in Chapter~22.

\section{Interpolation as an operator}
For fixed distinct nodes $x_0,\ldots,x_n$, define the interpolation operator
\[
(I_nf)(x)=\sum_{j=0}^n f(x_j)L_j(x),
\]
where $L_j$ are the Lagrange cardinal polynomials. Exact interpolation says $I_np=p$ for every polynomial $p$ of degree at most $n$. Stability asks a different question: how strongly can $I_n$ amplify perturbations in sampled values or in the target function?

\section{Lebesgue function and Lebesgue constant}
The \emph{Lebesgue function} and \emph{Lebesgue constant} of the nodes are
\[
\lambda_n(x)=\sum_{j=0}^n|L_j(x)|,
\qquad
\Lambda_n=\max_{x\in[a,b]}\lambda_n(x).
\]
They depend only on the node geometry, not on the sampled function.

\begin{theorem}[Lebesgue stability bound]\label{thm:ii20-lebesgue-stability}
For data perturbations $\delta_j$ at the nodes, the induced interpolation perturbation
\[
q(x)=\sum_{j=0}^n\delta_jL_j(x)
\]
satisfies
\[
\|q\|_\infty\le \Lambda_n\max_j|\delta_j|.
\]
In fact, the operator norm of $I_n:C[a,b]\to C[a,b]$ in the sup norm is exactly $\Lambda_n$.
\begin{proof}
For every $x$,
\[
|q(x)|\le\sum_j|\delta_j|\,|L_j(x)|
\le\lambda_n(x)\max_j|\delta_j|.
\]
Taking the maximum over $x$ gives the bound. For the reverse inequality, choose a point $x_*$ where $\lambda_n$ is maximal and choose a continuous function $g$ with $\|g\|_\infty\le1$ and nodal values $g(x_j)=\operatorname{sgn}(L_j(x_*))$; a piecewise-linear choice through the ordered nodes suffices. Then $|(I_ng)(x_*)|=\lambda_n(x_*)=\Lambda_n$.
\end{proof}
\end{theorem}

\begin{theorem}[Lebesgue quasi-optimality bound]\label{thm:ii20-lebesgue-quasioptimal}
Let
\[
E_n(f)=\inf_{p\in\mathcal P_n}\|f-p\|_\infty.
\]
Then
\[
\|f-I_nf\|_\infty\le(1+\Lambda_n)E_n(f).
\]
\begin{proof}
For any $p\in\mathcal P_n$, exact reproduction gives $I_np=p$. Hence
\[
f-I_nf=(f-p)-I_n(f-p),
\]
so
\[
\|f-I_nf\|_\infty
\le\|f-p\|_\infty+\|I_n(f-p)\|_\infty
\le(1+\Lambda_n)\|f-p\|_\infty.
\]
Taking the infimum over $p$ proves the claim.
\end{proof}
\end{theorem}

\begin{example}[Two-node interpolation has no sup-norm amplification]\label{ex:ii20-lebesgue-01}
For the endpoint nodes $-1,1$,
\[
L_0(x)=\frac{1-x}{2},\qquad L_1(x)=\frac{1+x}{2}.
\]
On $[-1,1]$ both are nonnegative and $L_0+L_1=1$, so
\[
\lambda_1(x)=1,\qquad \Lambda_1=1.
\]
Thus linear interpolation between two endpoints does not amplify the maximum nodal perturbation in the sup norm.
\end{example}

\section{Runge instability}
A small interpolation residual at the nodes does not imply small error between the nodes. The classical example is the analytic Runge function
\[
r(x)=\frac{1}{1+25x^2},\qquad -1\le x\le1.
\]
High-degree interpolation at equally spaced nodes can develop large endpoint oscillations even though every node is matched exactly. The phenomenon is not a contradiction of the Weierstrass theorem: polynomial \emph{approximation} can converge while this particular sequence of interpolation operators is unstable.

\begin{example}[Three different claims must not be confused]\label{ex:ii20-runge-01}
For equally spaced high-degree interpolation of the Runge function, all three statements can coexist:
\begin{enumerate}
\item the interpolating polynomial exists and is unique;
\item its residual at every interpolation node is exactly zero in exact arithmetic;
\item its sup-norm error between nodes can nevertheless be large.
\end{enumerate}
The first is algebra, the second is the interpolation constraint, and the third is a stability/approximation question governed in part by the growth of $\Lambda_n$.
\end{example}

\section{Why Chebyshev nodes help}
Chebyshev-type nodes cluster toward the endpoints and keep the interpolation operator much better conditioned in the sup norm. For the classical Chebyshev families the Lebesgue constant grows only logarithmically with degree, whereas for equally spaced nodes it grows exponentially. Chapter~22 develops the Chebyshev geometry from the minimax viewpoint.

\begin{remark}[Error formula versus stability factor]
The derivative remainder
\[
\frac{f^{(n+1)}(\xi)}{(n+1)!}\prod_j(x-x_j)
\]
and the Lebesgue constant answer complementary questions. The remainder formula can be sharp when high derivatives are controlled; the Lebesgue constant is function-independent and quantifies how node geometry amplifies data or approximation errors. Robust interpolation analysis should inspect both.
\end{remark}

% VOL02-NUMERICAL-INTERPOLATION-II20
\begin{remark}[Exact interpolation versus stable interpolation]
Distinct nodes guarantee a unique degree-at-most-\(n\) interpolating
polynomial in exact arithmetic.  This existence-and-uniqueness theorem does
not say that every coordinate representation of that polynomial is equally
well conditioned numerically.

In particular, solving for monomial coefficients through a Vandermonde system
can be sensitive when the nodes or the chosen polynomial basis produce a
large condition number.  Newton, Lagrange/barycentric, or orthogonal-polynomial
representations may describe the same exact interpolant while behaving very
differently in floating-point evaluation.

A small residual at the interpolation nodes,
\[
\max_j |p(x_j)-y_j|,
\]
checks agreement with the sampled data.  It does not by itself control the
error between nodes, nor does it certify that the coefficient vector is
well-conditioned.  Interpolation error, data sensitivity, and evaluation
stability are separate questions and should be reported separately when the
data are numerical.
\end{remark}


% BEGIN VOL02-EXPANSION II20-example-03
\begin{example}[Node choice controls the node polynomial]\label{ex:ii20-expand-03}
The interpolation error contains
\[
\omega_{n+1}(x)=\prod_{j=0}^n(x-x_j).
\]
For a fixed derivative bound, reducing
\(\|\omega_{n+1}\|_\infty\) reduces worst-case interpolation error. Chebyshev
nodes are designed precisely to control this factor on an interval. This
explains why node placement is not a cosmetic choice: equally spaced nodes can
produce a much larger oscillatory node polynomial near endpoints.
\end{example}
% END VOL02-EXPANSION II20-example-03
\section{Hermite preview}
Derivative data can be incorporated by repeated-node ideas.

\section{Core structural results}

\begin{theorem}[Interpolation existence and uniqueness]\label{thm:ii20-01}
For $n+1$ distinct nodes, there is a unique polynomial of degree at most $n$ taking prescribed values.
\begin{proof}
Lagrange basis polynomials give existence. The difference of two interpolants has $n+1$ roots but degree at most $n$, so it is zero.
\end{proof}
\end{theorem}

\begin{theorem}[Lagrange formula]\label{thm:ii20-02}
The interpolant is $p(x)=\sum_{j=0}^ny_jL_j(x)$ with $L_j(x_i)=\delta_{ij}$.
\begin{proof}
By construction each basis polynomial is one at its own node and zero at the others, so the sum matches all data.
\end{proof}
\end{theorem}

\begin{theorem}[Interpolation remainder]\label{thm:ii20-03}
Let $x_0,\ldots,x_n$ be distinct points in an interval $I$, let $p_n$ interpolate $f$ at those nodes, and assume $f\in C^{n+1}(I)$. For every $x\in I$ there is a point $\xi$ in the smallest interval containing $x,x_0,\ldots,x_n$ such that
\[
f(x)-p_n(x)=\frac{f^{(n+1)}(\xi)}{(n+1)!}\prod_{j=0}^n(x-x_j).
\]
\begin{proof}
Subtract a multiple of the node polynomial chosen to vanish also at $x$, then apply Rolle's theorem repeatedly to the resulting function.
\end{proof}
\end{theorem}

\section{Worked examples}

\begin{example}[Two-point interpolation]\label{ex:ii20-worked-01}
Interpolate $(0,1)$ and $(1,3)$. The unique linear interpolant is $p(x)=1+2x$.
\end{example}

\begin{example}[Three-point Lagrange]\label{ex:ii20-worked-02}
Interpolate $f(0)=0,f(1)=1,f(2)=4$. The data lie on $x^2$, so uniqueness gives $p(x)=x^2$.
\end{example}

\begin{example}[Uniqueness by roots]\label{ex:ii20-worked-03}
Why cannot two degree-$n$ interpolants agree at $n+1$ distinct nodes and differ elsewhere? Their difference would be a nonzero degree-at-most-$n$ polynomial with too many roots.
\end{example}

\begin{example}[Lagrange basis property]\label{ex:ii20-worked-04}
Verify $L_j(x_i)=\delta_{ij}$. At $i=j$ all factors equal one; for $i\ne j$ one numerator factor vanishes.
\end{example}

\section{Solved dossiers}
Each dossier is a substantial solved problem that combines several ideas from the chapter. The dossiers emphasize proof strategy, hypothesis checking, and connections between abstract results and concrete calculations.

\begin{problem}[Two-point interpolation]\label{prob:ii20-dossier-01}
Interpolate $(0,1)$ and $(1,3)$.
\end{problem}
\begin{solution}
The unique linear interpolant is $p(x)=1+2x$.
\end{solution}

\begin{problem}[Three-point Lagrange]\label{prob:ii20-dossier-02}
Interpolate $f(0)=0,f(1)=1,f(2)=4$.
\end{problem}
\begin{solution}
The data lie on $x^2$, so uniqueness gives $p(x)=x^2$.
\end{solution}

\begin{problem}[Uniqueness by roots]\label{prob:ii20-dossier-03}
Why cannot two degree-$n$ interpolants agree at $n+1$ distinct nodes and differ elsewhere?
\end{problem}
\begin{solution}
Their difference would be a nonzero degree-at-most-$n$ polynomial with too many roots.
\end{solution}

\begin{problem}[Lagrange basis property]\label{prob:ii20-dossier-04}
Verify $L_j(x_i)=\delta_{ij}$.
\end{problem}
\begin{solution}
At $i=j$ all factors equal one; for $i\ne j$ one numerator factor vanishes.
\end{solution}

\begin{problem}[Newton divided difference]\label{prob:ii20-dossier-05}
Compute the first divided difference for data $(x_0,y_0),(x_1,y_1)$.
\end{problem}
\begin{solution}
It is $(y_1-y_0)/(x_1-x_0)$.
\end{solution}

\begin{problem}[Incremental interpolation]\label{prob:ii20-dossier-06}
Why is Newton form convenient when adding one new node?
\end{problem}
\begin{solution}
The previous polynomial remains intact and one new product term is appended.
\end{solution}

\begin{problem}[Error node product]\label{prob:ii20-dossier-07}
Why do roots of the interpolation error include all nodes?
\end{problem}
\begin{solution}
The interpolant matches the function exactly at every node.
\end{solution}

\begin{problem}[Linear interpolation error]\label{prob:ii20-dossier-08}
State the error form for two nodes.
\end{problem}
\begin{solution}
$f(x)-p_1(x)=f''(\xi)(x-x_0)(x-x_1)/2$.
\end{solution}

\begin{problem}[Runge warning]\label{prob:ii20-dossier-09}
What can go wrong with high-degree equally spaced interpolation?
\end{problem}
\begin{solution}
Endpoint oscillations can grow despite exact matching at nodes.
\end{solution}

\begin{problem}[Chebyshev motivation]\label{prob:ii20-dossier-10}
Why choose nodes near Chebyshev points?
\end{problem}
\begin{solution}
They reduce the maximum magnitude of the node polynomial and thereby control worst-case interpolation error.
\end{solution}

\begin{problem}[Vandermonde formulation]\label{prob:ii20-dossier-11}
How can interpolation be written as a linear system?
\end{problem}
\begin{solution}
Writing $p(x)=\sum a_kx^k$ gives a Vandermonde matrix whose invertibility follows from distinct nodes.
\end{solution}

\begin{problem}[Interpolation synthesis]\label{prob:ii20-dossier-12}
What is being approximated and what is exact?
\end{problem}
\begin{solution}
The polynomial matches finitely many data exactly, while the remainder formula quantifies error between those nodes in terms of higher derivatives and node geometry.
\end{solution}

\begin{problem}[Lebesgue amplification]\label{prob:ii20-dossier-13}
If the nodal data are perturbed by at most $\varepsilon$, what worst-case sup-norm perturbation does the Lebesgue constant certify?
\end{problem}
\begin{solution}
At most $\Lambda_n\varepsilon$.
\end{solution}

\begin{problem}[Runge phenomenon versus Weierstrass]\label{prob:ii20-dossier-14}
Why does divergence or severe oscillation of equally spaced interpolants for the Runge function not contradict polynomial density in $C[-1,1]$?
\end{problem}
\begin{solution}
Weierstrass asserts that some sequence of polynomials approximates every continuous function uniformly. It does not assert that the interpolation polynomial on every prescribed node family converges. Equally spaced interpolation can be an unstable choice of approximation operator.
\end{solution}

% VOL02-HINT-RECONCILIATION-II20
\section{Exercises with complete solutions}

\begin{exercise}\label{exr:ii20-01}
Interpolate two points $(0,0),(1,1)$.
\end{exercise}
\begin{hint}
Line through them.
\end{hint}
\begin{solution}
$p=x$.
\end{solution}

\begin{exercise}\label{exr:ii20-02}
How many data points determine degree at most $n$?
\end{exercise}
\begin{hint}
A degree-at-most-n polynomial is uniquely determined by n+1 values at distinct nodes.
\end{hint}
\begin{solution}
$n+1$.
\end{solution}

\begin{exercise}\label{exr:ii20-03}
Why distinct nodes?
\end{exercise}
\begin{hint}
Vandermonde determinant.
\end{hint}
\begin{solution}
Otherwise uniqueness fails or data may conflict.
\end{solution}

\begin{exercise}\label{exr:ii20-04}
What are Lagrange basis values at nodes?
\end{exercise}
\begin{hint}
Kronecker delta.
\end{hint}
\begin{solution}
$0$ or $1$.
\end{solution}

\begin{exercise}\label{exr:ii20-05}
What degree is the node polynomial?
\end{exercise}
\begin{hint}
One factor per node.
\end{hint}
\begin{solution}
$n+1$.
\end{solution}

\begin{exercise}\label{exr:ii20-06}
What theorem proves uniqueness?
\end{exercise}
\begin{hint}
Subtract two candidates; their difference would have n+1 distinct zeros but degree at most n.
\end{hint}
\begin{solution}
A nonzero degree-$n$ polynomial has at most $n$ roots.
\end{solution}

\begin{exercise}\label{exr:ii20-07}
Does interpolation guarantee small error?
\end{exercise}
\begin{hint}
Interpolation is exact at the nodes but may oscillate badly between them; inspect the remainder and the node placement.
\end{hint}
\begin{solution}
Node placement and smoothness matter.
\end{solution}

\begin{exercise}\label{exr:ii20-08}
What form supports incremental data?
\end{exercise}
\begin{hint}
Newton form.
\end{hint}
\begin{solution}
Newton divided-difference form.
\end{solution}


% BEGIN VOL02-EXPANSION II20-exercises-01
\section{Graded supplementary exercises}
These exercises extend the core exercise set and are grouped by mathematical role and increasing level of synthesis.

\subsection*{Standard computations}

\begin{exercise}[Linear interpolation]\label{exr:ii20-expand-01}
Interpolate \((0,2)\) and \((2,6)\).
\end{exercise}
\begin{hint}
Find the line through the points.
\end{hint}
\begin{solution}
\(p(x)=2+2x\).
\end{solution}

\begin{exercise}[Lagrange basis]\label{exr:ii20-expand-02}
For nodes \(0,1,2\), write \(L_1(x)\).
\end{exercise}
\begin{hint}
Use products over the other nodes.
\end{hint}
\begin{solution}
\(L_1(x)=\frac{x(x-2)}{(1-0)(1-2)}=-x(x-2)\).
\end{solution}

\begin{exercise}[Divided difference]\label{exr:ii20-expand-03}
Compute \(f[0,1]\) for values \(f(0)=1,f(1)=4\).
\end{exercise}
\begin{hint}
Use the first divided-difference formula.
\end{hint}
\begin{solution}
\(f[0,1]=3\).
\end{solution}

\begin{exercise}[Vandermonde determinant]\label{exr:ii20-expand-04}
What condition on nodes makes the Vandermonde matrix invertible?
\end{exercise}
\begin{hint}
Recall its determinant product.
\end{hint}
\begin{solution}
The nodes must be pairwise distinct.
\end{solution}

\begin{exercise}[Remainder factor]\label{exr:ii20-expand-05}
For quadratic interpolation at nodes \(0,1,2\), what node product appears in the remainder?
\end{exercise}
\begin{hint}
Multiply \(x-x_j\).
\end{hint}
\begin{solution}
\(x(x-1)(x-2)\).
\end{solution}

\subsection*{Proofs}

\begin{exercise}[Uniqueness]\label{exr:ii20-expand-06}
Prove a degree-at-most-\(n\) interpolant at \(n+1\) distinct nodes is unique.
\end{exercise}
\begin{hint}
Subtract two interpolants.
\end{hint}
\begin{solution}
The difference has \(n+1\) distinct roots but degree at most \(n\), so it is zero.
\end{solution}

\begin{exercise}[Lagrange existence]\label{exr:ii20-expand-07}
Prove the Lagrange formula matches prescribed node values.
\end{exercise}
\begin{hint}
Use \(L_j(x_i)=\delta_{ij}\).
\end{hint}
\begin{solution}
At node \(x_i\), every term vanishes except \(y_iL_i(x_i)=y_i\).
\end{solution}

\begin{exercise}[Newton update]\label{exr:ii20-expand-08}
Explain why adding one node adds only one term to Newton form.
\end{exercise}
\begin{hint}
The existing interpolant already matches old nodes; multiply the correction by the old node polynomial.
\end{hint}
\begin{solution}
The correction \(c\prod_{j=0}^n(x-x_j)\) vanishes at all old nodes, and \(c\) is chosen to fit the new one.
\end{solution}

\begin{exercise}[Remainder theorem mechanism]\label{exr:ii20-expand-09}
Outline the repeated Rolle proof of the interpolation remainder.
\end{exercise}
\begin{hint}
Subtract a multiple of the node polynomial so the error also vanishes at the evaluation point.
\end{hint}
\begin{solution}
The constructed function has \(n+2\) zeros; repeated Rolle gives a zero of its \((n+1)\)-st derivative and yields the formula.
\end{solution}

\subsection*{Counterexamples and hypothesis testing}

\begin{exercise}[Repeated nodes]\label{exr:ii20-expand-10}
Why do repeated nodes break ordinary Lagrange interpolation formulas?
\end{exercise}
\begin{hint}
Denominators contain node differences.
\end{hint}
\begin{solution}
Distinctness fails and denominators vanish; derivative data require Hermite-type interpolation instead.
\end{solution}

\begin{exercise}[High-degree instability]\label{exr:ii20-expand-11}
Why can equally spaced high-degree interpolation behave poorly near interval endpoints?
\end{exercise}
\begin{hint}
Node-product and Lebesgue factors can grow strongly.
\end{hint}
\begin{solution}
Runge-type oscillations can amplify data/approximation error despite exact interpolation at nodes.
\end{solution}

\begin{exercise}[Interpolation not approximation optimum]\label{exr:ii20-expand-12}
Why need an interpolating polynomial not be the best uniform approximant of the same degree?
\end{exercise}
\begin{hint}
Interpolation enforces exact values at chosen nodes, while minimax optimizes worst-case error globally.
\end{hint}
\begin{solution}
The two optimization criteria are different.
\end{solution}

\subsection*{Applications and investigations}

\begin{exercise}[Table update]\label{exr:ii20-expand-13}
Why is Newton form attractive for streaming data?
\end{exercise}
\begin{hint}
New data append divided differences without changing old coefficients.
\end{hint}
\begin{solution}
It supports incremental interpolation with less recomputation.
\end{solution}

\begin{exercise}[Error budgeting]\label{exr:ii20-expand-14}
How can derivative and node-product bounds guide node count?
\end{exercise}
\begin{hint}
Use the remainder inequality.
\end{hint}
\begin{solution}
Choose degree/nodes until \(M\|\omega\|_\infty/(n+1)!\) falls below the desired tolerance.
\end{solution}

\subsection*{Challenge problems}

\begin{exercise}[Chebyshev node motivation]\label{exr:ii20-expand-15}
Explain why roots of a Chebyshev polynomial are natural interpolation nodes.
\end{exercise}
\begin{hint}
They minimize the size of an associated monic node polynomial up to scaling.
\end{hint}
\begin{solution}
This reduces the worst-case factor in the interpolation remainder.
\end{solution}

\begin{exercise}[Hermite dimension count]\label{exr:ii20-expand-16}
Why do values and first derivatives at \(m\) distinct nodes determine a polynomial of degree at most \(2m-1\) uniquely?
\end{exercise}
\begin{hint}
The difference of two candidates has a double zero at each node.
\end{hint}
\begin{solution}
It would have at least \(2m\) zeros counting multiplicity but degree at most \(2m-1\), so it must vanish.
\end{solution}
\begin{exercise}[Lebesgue function]\label{exr:ii20-lebesgue-01}
For nodes $-1,1$, compute the Lebesgue constant.
\end{exercise}
\begin{hint}
The two Lagrange basis polynomials are nonnegative and sum to one on the interval.
\end{hint}
\begin{solution}
$\Lambda_1=1$.
\end{solution}

\begin{exercise}[Perturbed samples]\label{exr:ii20-lebesgue-02}
If $\Lambda_n=8$ and every sampled value has absolute error at most $10^{-5}$, bound the interpolation error caused by the data perturbation alone.
\end{exercise}
\begin{hint}
Apply the Lebesgue stability bound.
\end{hint}
\begin{solution}
The induced sup-norm perturbation is at most $8\times10^{-5}$.
\end{solution}

\begin{exercise}[Quasi-optimality]\label{exr:ii20-lebesgue-03}
If $E_n(f)=10^{-4}$ and $\Lambda_n=5$, what bound follows for $\|f-I_nf\|_\infty$?
\end{exercise}
\begin{hint}
Use $(1+\Lambda_n)E_n(f)$.
\end{hint}
\begin{solution}
At most $6\times10^{-4}$.
\end{solution}

\begin{exercise}[Runge diagnosis]\label{exr:ii20-runge-01}
A high-degree interpolant has zero residual at every node but large oscillation near the endpoints. Which diagnostic from this section directly measures sensitivity to the node geometry?
\end{exercise}
\begin{hint}
Inspect the operator norm of the interpolation map.
\end{hint}
\begin{solution}
The Lebesgue constant $\Lambda_n$.
\end{solution}

% END VOL02-EXPANSION II20-exercises-01
\section*{Chapter summary}
The chapter's tools form part of the Volume II analysis toolkit and will be used in later approximation, fixed-point, and convergence arguments.
